% WHERE A FENCED FRACTION IS SPLIT.
%
% `\left(AB\over C\right)' is ONE fraction standing between the two
% delimiters, and its numerator is what was read before the \over -- the
% opening delimiter LEFT OUT of it:
%
%   \hbox                          the fence
%   .\hbox(0.0+0.0)x0.0            the \left delimiter
%   .\hbox                         the fraction, delimiters and all
%   ..\hbox(0.0+0.0)x0.0           its own \nulldelimiterspace
%   ..\vbox
%   ...\hbox   ( A B               the numerator
%   ...\kern \rule \kern
%   ...\hbox   C                   the denominator, widened to match
%   ..\hbox(0.0+0.0)x0.0
%   .\hbox(0.0+0.0)x0.0            the \right delimiter
%
% (The `(' and `)' stand in the list as ORDINARY characters, since INITEX
% gives them no \delcode and the reference puts a null delimiter in their
% place.)
%
% This engine put an opening and a closing atom into the list before it
% set it -- they earn the right spacing when the group is not a fraction --
% and they moved the split by one, so the numerator came out `( A' and the
% denominator `B C'. An opening before an inner atom and an inner before a
% closing are spaced by nothing, so a fenced fraction needs neither.
%
% trip line 273 writes \left(...\over...\right and line 283 another.
\catcode`\{=1 \catcode`\}=2 \catcode`\$=3
\tracingonline=1 \showboxdepth=9 \showboxbreadth=99
\font\r=cmr10 \font\s=cmsy10 \font\x=cmex10 \r
\textfont0=\r \scriptfont0=\r \scriptscriptfont0=\r
\textfont1=\r \scriptfont1=\r \scriptscriptfont1=\r
\textfont2=\s \scriptfont2=\s \scriptscriptfont2=\s
\textfont3=\x \scriptfont3=\x \scriptscriptfont3=\x
\message{[A an ordinary fraction]}
\setbox0=\hbox{$AB\over C$}
\showbox0
\message{[B one inside \left ... \right]}
\setbox0=\hbox{$\left(AB\over C\right)$}
\showbox0
\message{[C one whose numerator is only the delimiter]}
\setbox0=\hbox{$\left(\over C\right)$}
\showbox0
\message{[done]}
\end
